
\section{滤过 $A_\infty$ 代数情形}
\label{sec:8-5}

回忆，在研究与拉格朗日子流形 $L$ 关联的滤过 $A_\infty$ 代数结构时，
我们使用模空间的主分支。（主分支的定义见第~\ref{sec:2-1-2} 小节。）
因此，我们选取边界标记点的次序，使之与边界的逆时针定向相容。设
$[P_i,f_i]\in C^{g_i}(L;\Lambda_{0,\mathrm{nov}})$。由定义
\ref{def:8-4-1}，有
\[
\begin{aligned}
 &\mathcal M_{k+1}^{\mathrm{main}}(\beta;P_1,\ldots,P_k)\\
 &\quad:={}
 (-1)^{\epsilon_1}
 \mathcal M_{k+1}^{\mathrm{main}}(L;\beta)_
 {(\operatorname{ev}_1,\ldots,\operatorname{ev}_k)}
 \mathop{\times}_{f_1\times\cdots\times f_k}
 \left(\prod_{i=1}^{k}P_i\right),
\end{aligned}
\]
其中
\[
 \epsilon_1=(n+1)
 \sum_{j=1}^{k-1}\sum_{i=1}^{j}\deg P_i.
\]

我们来核对符号，从而完成定理~\ref{thm:3-5-11} 的证明；该定理断言
$C(L;\Lambda_{0,\mathrm{nov}})$ 连同运算 $\mathfrak m$ 构成滤过
$A_\infty$ 代数。
需要核对符号之处有两个：第一处在
$\widehat d\circ\widehat d=0$ 的证明中；第二处是比较我们在下列对象上的
定向与 $P_1\cap P_2$ 上的几何定向：
\[
 \bigl(\mathcal M_3^{\mathrm{main}}(L;\beta_0;P_1,P_2),
 \operatorname{ev}_0\bigr).
\]
（注意 $\beta_0=0$）。本节核对第一处，下一节核对第二处。为了分别证明
\eqref{eq:3-5-17-1} 和 \eqref{eq:3-5-17-2}，我们分别需要下面命题
\ref{prop:8-5-1} 的 (2) 和 (1)。

\newpage
\begin{proposition}{8.5.1}\label{prop:8-5-1}
\begin{enumerate}[label=\textup{(\arabic*)},leftmargin=2.8em]
\item 对 $\beta=\beta_1+\beta_2$，有
\[
\begin{aligned}
 &\mathcal M_{k-k_2+2}^{\mathrm{main}}
 \bigl(\beta_1;P_1,\ldots,P_{i-1},
 \mathcal M_{k_2+1}^{\mathrm{main}}
 (\beta_2;P_i,\ldots,P_{i+k_2-1}),
 P_{i+k_2},\ldots,P_k\bigr)\\
 &\quad\subset
 (-1)^{\epsilon_2}\partial
 \mathcal M_{k+1}^{\mathrm{main}}
 \bigl(\beta;P_1,\ldots,P_{i-1},P_i,\ldots,
 P_{i+k_2-1},P_{i+k_2},\ldots,P_k\bigr),
\end{aligned}
\]
其中
\[
 \epsilon_2=n+1+\sum_{j=1}^{i-1}(\deg P_j+1).
\]

\item 有
\[
\begin{aligned}
 &\mathcal M_{k+1}^{\mathrm{main}}
 (\beta;P_1,\ldots,\partial P_i,\ldots,P_k)\\
 &\quad\subset
 (-1)^{\epsilon_3}\partial
 \mathcal M_{k+1}^{\mathrm{main}}
 (\beta;P_1,\ldots,P_i,\ldots,P_k),
\end{aligned}
\]
其中
\[
 \epsilon_3=1+\sum_{j=1}^{i-1}(\deg P_j+1).
\]
\end{enumerate}
\end{proposition}

\begin{proof}
(1) 的证明分为若干步。

\par\smallskip\noindent\textbf{第 1 步。}
注意，除循环置换外，改变标记点次序并不保持主分支。不过，在计算的
中间阶段，我们将使用主分支以外的分支。由命题 3.5.2，注意到
\[
 \deg\mathcal M_{k_2+1}^{\mathrm{main}}
 (\beta_2;P_i,\ldots,P_{i+k_2-1})
 \equiv\sum_{j=i}^{i+k_2-1}(\deg P_j+1)\pmod2.
\]
反复使用引理~\ref{lem:8-4-3}，得到
\[
\begin{aligned}
 &\mathcal M_{k-k_2+2}^{\mathrm{main}}
 \bigl(\beta_1;P_1,\ldots,
 \mathcal M_{k_2+1}^{\mathrm{main}}
 (\beta_2;P_i,\ldots,P_{i+k_2-1}),
 P_{i+k_2},\ldots,P_k\bigr)\\
 &\quad\subset
 (-1)^{\delta_1}\mathcal M_{k-k_2+2}
 \bigl(\beta_1;
 \mathcal M_{k_2+1}^{\mathrm{main}}
 (\beta_2;P_i,\ldots,P_{i+k_2-1}),\\[-2pt]
 &\hspace{8em}P_1,\ldots,P_{i-1},
 P_{i+k_2},\ldots,P_k\bigr),
\end{aligned}
\]
其中
\[
 \delta_1=\left(1+\sum_{j=i}^{i+k_2-1}(\deg P_j+1)\right)
 \left(\sum_{j=1}^{i-1}(\deg P_j+1)\right).
\]

\par\smallskip\noindent\textbf{第 2 步。}
下面比较
\[
\begin{aligned}
 \mathcal M_{k-k_2+2}\bigl(&\beta_1;
 \mathcal M_{k_2+1}^{\mathrm{main}}
 (\beta_2;P_i,\ldots,P_{i+k_2-1}),\\[-2pt]
 &P_1,\ldots,P_{i-1},P_{i+k_2},\ldots,P_k\bigr)
\end{aligned}
\]
上的定向与
\[
 \partial\mathcal M_{k+1}
 (\beta;P_i,\ldots,P_{i+k_2-1},P_1,\ldots,
 P_{i-1},P_{i+k_2},\ldots,P_k)
\]
上的定向。由定义~\ref{def:8-4-1}，有
\[
\begin{aligned}
 &\mathcal M_{k+1}
 (\beta;P_i,\ldots,P_{i+k_2-1},P_1,\ldots,
 P_{i-1},P_{i+k_2},\ldots,P_k)\\
 &\quad={}
 (-1)^{\gamma_1}\mathcal M_{k+1}(\beta)_
 {(\operatorname{ev}_1,\ldots,\operatorname{ev}_k)}
 \times\left(
 \prod_{j=i}^{i+k_2-1}P_j\times
 \prod_{j=1}^{i-1}P_j\times
 \prod_{j=i+k_2}^{k}P_j\right),
\end{aligned}
\]
其中
\[
\begin{aligned}
 \gamma_1=(n+1)\Bigg(&
 \sum_{j=i}^{i+k_2-1}\sum_{\ell=i}^{j}\deg P_\ell\\
 &+\sum_{j=1}^{i-1}\left(
   \sum_{m=i}^{i+k_2-1}\deg P_m
   +\sum_{\ell=1}^{j}\deg P_\ell\right)\\
 &+\sum_{j=i+k_2}^{k-1}\left(
   \sum_{m=i}^{i+k_2-1}\deg P_m
   +\sum_{m=1}^{i-1}\deg P_m
   +\sum_{\ell=i+k_2}^{j}\deg P_\ell\right)\Bigg).
\end{aligned}
\]
使用引理~\ref{lem:8-2-3}(3) 的迭代公式，得到
\[
\begin{aligned}
 &\mathcal M_{k+1}(\beta)_
 {(\operatorname{ev}_1,\ldots,\operatorname{ev}_k)}
 \times\left(
 \prod_{j=i}^{i+k_2-1}P_j\times
 \prod_{j=1}^{i-1}P_j\times
 \prod_{j=i+k_2}^{k}P_j\right)\\
 &\quad={}
 (-1)^{\gamma_2}
 \left(
 \mathcal M_{k+1}(\beta)_
 {(\operatorname{ev}_1,\ldots,\operatorname{ev}_{k_2})}
 \times\prod_{j=i}^{i+k_2-1}P_j
 \right)_
 {(\operatorname{ev}_{k_2+1},\ldots,\operatorname{ev}_k)}\\[-2pt]
 &\hspace{13em}{}\times
 \left(\prod_{j=1}^{i-1}P_j\times
 \prod_{j=i+k_2}^{k}P_j\right),
\end{aligned}
\]
其中
\[
\begin{aligned}
 \gamma_2
 &=(k-k_2)n\left(k_2n+
 \sum_{j=i}^{i+k_2-1}(n-\deg P_j)\right)\\
 &\equiv(k-k_2)n\sum_{j=i}^{i+k_2-1}\deg P_j\pmod2.
\end{aligned}
\]

粘合两个带标记点的全纯圆盘时，我们关于标记点次序与边界定向的约定
（见注~\ref{rem:8-3-4}）以及命题~\ref{prop:8-3-3} 表明
\[
\begin{aligned}
 &\partial\Bigg[
 \left(
 \mathcal M_{k+1}(\beta)_
 {(\operatorname{ev}_1,\ldots,\operatorname{ev}_{k_2})}
 \times\prod_{j=i}^{i+k_2-1}P_j
 \right)_
 {(\operatorname{ev}_{k_2+1},\ldots,\operatorname{ev}_k)}\\[-2pt]
 &\hspace{10em}{}\times
 \left(\prod_{j=1}^{i-1}P_j\times
 \prod_{j=i+k_2}^{k}P_j\right)\Bigg]\\
 &\quad\supset(-1)^{\gamma_3}
 \Bigg[
 \begin{gathered}
 \left(
 \mathcal M_{k-k_2+2}(\beta_1)_
 {\operatorname{ev}_1^{\beta_1}}
 \mathop{\times}_{\operatorname{ev}_0^{\beta_2}}
 \mathcal M_{k_2+1}(\beta_2)
 \right)_
 {(\operatorname{ev}_1^\beta,\ldots,
   \operatorname{ev}_{k_2}^\beta)}\\[-2pt]
 {}\times\displaystyle\prod_{j=i}^{i+k_2-1}P_j
 \end{gathered}
 \Bigg]_
 {(\operatorname{ev}_{k_2+1}^\beta,\ldots,
   \operatorname{ev}_{k}^\beta)}\\[-2pt]
 &\hspace{10em}{}\times
 \left(\prod_{j=1}^{i-1}P_j\right)
 \times\left(\prod_{j=i+k_2}^{k}P_j\right),
\end{aligned}
\]
其中
\[
\begin{aligned}
 \gamma_3
 &=(k-k_2)(k_2-1)+n+k-k_2+1\\
 &\equiv n+1+kk_2+k_2\pmod2.
\end{aligned}
\]
由引理~\ref{lem:8-2-3}(2) 的结合性质与定义~\ref{def:8-4-1}，有
\begingroup
\small
\[
\begin{aligned}
 &\Bigg[
 \left(
 \mathcal M_{k-k_2+2}(\beta_1)_
 {\operatorname{ev}_1^{\beta_1}}
 \mathop{\times}_{\operatorname{ev}_0^{\beta_2}}
 \mathcal M_{k_2+1}(\beta_2)
 \right)_
 {(\operatorname{ev}_1^\beta,\ldots,
   \operatorname{ev}_{k_2}^\beta)}
 \times\prod_{j=i}^{i+k_2-1}P_j
 \Bigg]_
 {(\operatorname{ev}_{k_2+1}^\beta,\ldots,
   \operatorname{ev}_{k}^\beta)}\\[-2pt]
 &\hspace{11em}{}\times
 \left(\prod_{j=1}^{i-1}P_j\times
 \prod_{j=i+k_2}^{k}P_j\right)\\
 &={}
 \Bigg[
 \mathcal M_{k-k_2+2}(\beta_1)_
 {\operatorname{ev}_1^{\beta_1}}
 \mathop{\times}_{\operatorname{ev}_0^{\beta_2}}
 \left(
 \mathcal M_{k_2+1}(\beta_2)_
 {(\operatorname{ev}_1^{\beta_2},\ldots,
   \operatorname{ev}_{k_2}^{\beta_2})}
 \times\prod_{j=i}^{i+k_2-1}P_j
 \right)
 \Bigg]_
 {(\operatorname{ev}_{k_2+1}^\beta,\ldots,
   \operatorname{ev}_{k}^\beta)}\\[-2pt]
 &\hspace{11em}{}\times
 \left(\prod_{j=1}^{i-1}P_j\times
 \prod_{j=i+k_2}^{k}P_j\right)\\
 &={}(-1)^{\gamma_4}\\[-2pt]
 &\left(
 \mathcal M_{k-k_2+2}(\beta_1)_
 {\operatorname{ev}_1^{\beta_1}}
 \!\mathop{\times}_{\operatorname{ev}_0^{\beta_2}}\!
 \mathcal M_{k_2+1}
 (\beta_2;P_i,\ldots,P_{i+k_2-1})
 \right)_
 {(\operatorname{ev}_{k_2+1}^\beta,\ldots,
   \operatorname{ev}_{k}^\beta)}\\[-2pt]
 &\hspace{11em}{}\times
 \left(\prod_{j=1}^{i-1}P_j\times
 \prod_{j=i+k_2}^{k}P_j\right),
\end{aligned}
\]
\endgroup
其中
\[
 \gamma_4=(n+1)
 \sum_{j=i}^{i+k_2-2}\sum_{\ell=i}^{j}\deg P_\ell.
\]

再次使用迭代公式，得到
\[
\begin{aligned}
 &\left(
 \mathcal M_{k-k_2+2}(\beta_1)_
 {\operatorname{ev}_1^{\beta_1}}
 \mathop{\times}_{\operatorname{ev}_0^{\beta_2}}
 \mathcal M_{k_2+1}
 (\beta_2;P_i,\ldots,P_{i+k_2-1})
 \right)_
 {(\operatorname{ev}_{k_2+1}^\beta,\ldots,
   \operatorname{ev}_{k}^\beta)}\\[-2pt]
 &\hspace{11em}{}\times
 \left(\prod_{j=1}^{i-1}P_j\times
 \prod_{j=i+k_2}^{k}P_j\right)\\
 &={}
 (-1)^{\gamma_5}
 \mathcal M_{k-k_2+2}(\beta_1)_
 {(\operatorname{ev}_1^{\beta_1},\ldots,
   \operatorname{ev}_{k-k_2+1}^{\beta_1})}
 \times\Bigg(
 \mathcal M_{k_2+1}(\beta_2;P_i,\ldots,P_{i+k_2-1})\\[-2pt]
 &\hspace{18em}{}\times
 \left(\prod_{j=1}^{i-1}P_j\times
 \prod_{j=i+k_2}^{k}P_j\right)\Bigg)\\
 &={}
 (-1)^{\gamma_5+\gamma_6}
 \mathcal M_{k-k_2+2}
 \bigl(\beta_1;
 \mathcal M_{k_2+1}(\beta_2;P_i,\ldots,P_{i+k_2-1}),\\[-2pt]
 &\hspace{8em}P_1,\ldots,P_{i-1},
 P_{i+k_2},\ldots,P_k\bigr),
\end{aligned}
\]
其中
\[
\begin{aligned}
 \gamma_5
 &=n(k-k_2)\left(2n-
   \sum_{j=i}^{i+k_2-1}(\deg P_j+1)\right)\\
 &\equiv n(k-k_2)\left(k_2+
   \sum_{j=i}^{i+k_2-1}\deg P_j\right)\pmod2,
\end{aligned}
\]
以及
\[
\begin{aligned}
 \gamma_6
 ={}&(n+1)\Bigg(
 \deg\mathcal M_{k_2+1}(\beta_2;P_i,\ldots,P_{i+k_2-1})\\
 &+\sum_{j=1}^{i-1}\left(
   \deg\mathcal M_{k_2+1}(\beta_2;P_i,\ldots,P_{i+k_2-1})\right.\\[-2pt]
 &\hspace{11em}\left.{}+
   \sum_{\ell=1}^{j}\deg P_\ell\right)\\
 &+\sum_{j=i+k_2}^{k-1}\left(
   \deg\mathcal M_{k_2+1}(\beta_2;P_i,\ldots,P_{i+k_2-1})\right.\\[-2pt]
 &\hspace{11em}\left.{}+
   \sum_{m=1}^{i-1}\deg P_m
   +\sum_{\ell=i+k_2}^{j}\deg P_\ell\right)\Bigg)\\
 &=(n+1)\Bigg(
 (k-k_2)\left(k_2+
   \sum_{m=i}^{i+k_2-1}\deg P_m\right)\\
 &\hspace{5em}{}
 +(k-i-k_2)\left(\sum_{m=1}^{i-1}\deg P_m\right)\\
 &+\sum_{j=1}^{i-1}\sum_{\ell=1}^{j}\deg P_\ell
 +\sum_{j=i+k_2}^{k-1}\sum_{\ell=i+k_2}^{j}\deg P_\ell
 \Bigg).
\end{aligned}
\]
于是，初等计算给出
\[
 \gamma_1+\gamma_2+\gamma_3+\gamma_4+\gamma_5+\gamma_6
 \equiv n+1\pmod2.
\]
因而我们得到以下包含关系；其中外层圆盘的同伦类与外层模空间的下标
均按这一分解的前后等式统一。%
\translatornote{源书在上一串等式的最后一行把外层圆盘类
$\beta_1$ 误写成总类 $\beta$；又在下式把
$\mathcal M_{k-k_2+2}$ 误印成 $\mathcal M_{k-k_2+1}$。
命题~\ref{prop:8-5-1}(1)、本步此前各式，以及外层共有
$k-k_2+1$ 个输入，分别唯一确定这里应为 $\beta_1$ 与下标
$k-k_2+2$。}
\[
\begin{aligned}
 &\mathcal M_{k-k_2+2}
 \bigl(\beta_1;
 \mathcal M_{k_2+1}^{\mathrm{main}}
 (\beta_2;P_i,\ldots,P_{i+k_2-1}),
 P_1,\ldots,P_{i-1},P_{i+k_2},\ldots,P_k\bigr)\\
 &\quad\subset
 (-1)^{n+1}\partial\mathcal M_{k+1}
 (\beta;P_i,\ldots,P_{i+k_2-1},P_1,\ldots,
 P_{i-1},P_{i+k_2},\ldots,P_k).
\end{aligned}
\]

\par\smallskip\noindent\textbf{第 3 步。}
另一方面，再次使用引理~\ref{lem:8-4-3} 可见
\[
\begin{aligned}
 &\mathcal M_{k+1}
 (\beta;P_i,\ldots,P_{i+k_2-1},P_1,\ldots,
 P_{i-1},P_{i+k_2},\ldots,P_k)\\
 &\quad={}
 (-1)^{\delta_2}\mathcal M_{k+1}^{\mathrm{main}}
 (\beta;P_1,\ldots,P_{i-1},P_i,\ldots,
 P_{i+k_2-1},P_{i+k_2},\ldots,P_k),
\end{aligned}
\]
其中
\[
 \delta_2=\left(\sum_{j=i}^{i+k_2-1}(\deg P_j+1)\right)
 \left(\sum_{j=1}^{i-1}(\deg P_j+1)\right).
\]
因此
\[
 \delta_1+n+1+\delta_2
 \equiv n+1+\sum_{j=1}^{i-1}(\deg P_j+1)\pmod2,
\]
这就证明了命题~\ref{prop:8-5-1}(1)。

\par\smallskip
(2) 现在证明命题~\ref{prop:8-5-1}(2)。回忆
\[
\begin{aligned}
 &\mathcal M_{k+1}^{\mathrm{main}}
 (\beta;P_1,\ldots,P_i,\ldots,P_k)\\
 &\quad={}
 (-1)^{\epsilon_1}
 \mathcal M_{k+1}^{\mathrm{main}}(\beta)_
 {(\operatorname{ev}_1,\ldots,\operatorname{ev}_k)}
 \mathop{\times}_{f_1\times\cdots\times f_k}
 \left(\prod_{j=1}^{k}P_j\right),
\end{aligned}
\]
其中
\[
 \epsilon_1=(n+1)
 \sum_{j=1}^{k-1}\sum_{\ell=1}^{j}\deg P_\ell.
\]
由引理~\ref{lem:8-2-3}(1)，有
\[
\begin{aligned}
 &\partial\left(
 \mathcal M_{k+1}^{\mathrm{main}}(\beta)_
 {(\operatorname{ev}_1,\ldots,\operatorname{ev}_k)}
 \mathop{\times}_{f_1\times\cdots\times f_k}
 \left(\prod_{j=1}^{k}P_j\right)\right)\\
 &={}
 \partial\mathcal M_{k+1}^{\mathrm{main}}(\beta)_
 {(\operatorname{ev}_1,\ldots,\operatorname{ev}_k)}
 \mathop{\times}_{f_1\times\cdots\times f_k}
 \left(\prod_{j=1}^{k}P_j\right)\\
 &\quad\sqcup
 (-1)^{n+k+nk}
 \mathcal M_{k+1}^{\mathrm{main}}(\beta)_
 {(\operatorname{ev}_1,\ldots,\operatorname{ev}_k)}
 \mathop{\times}_{f_1\times\cdots\times f_k}
 \partial\left(\prod_{j=1}^{k}P_j\right),
\end{aligned}
\]
因为
\[
 \dim\mathcal M_{k+1}^{\mathrm{main}}(\beta)
 =n+\mu(\beta)-3+k+1\equiv n+k\pmod2.
\]
此外，再次使用引理~\ref{lem:8-2-3}，容易看出
\[
 \partial\left(\prod_{j=1}^{k}P_j\right)
 =\coprod_{i=1}^{k}
 (-1)^{\sum_{j=1}^{i-1}\dim P_j}
 P_1\times\cdots\times\partial P_i\times\cdots\times P_k.
\]
另一方面，定义~\ref{def:8-4-1} 给出
\[
\begin{aligned}
 &\mathcal M_{k+1}^{\mathrm{main}}
 (\beta;P_1,\ldots,\partial P_i,\ldots,P_k)\\
 &\quad={}
 (-1)^{\epsilon_2}
 \mathcal M_{k+1}^{\mathrm{main}}(\beta)_
 {(\operatorname{ev}_1,\ldots,\operatorname{ev}_k)}
 \mathop{\times}_{f_1\times\cdots\times f_k}
 \left(\prod_{j=1}^{i-1}P_j\times\partial P_i\times
 \prod_{j=i+1}^{k}P_j\right),
\end{aligned}
\]
其中
\[
 \epsilon_2=(n+1)
 \sum_{j=1}^{k-1}\sum_{\ell=1}^{j}(\deg P_\ell)',
\]
而
\[
 (\deg P_\ell)'=
 \begin{cases}
  \deg P_\ell,&\ell\ne i,\\
  \deg P_\ell+1,&\ell=i.
 \end{cases}
\]
因此可见
\[
\begin{aligned}
 &\epsilon_1+n+k+nk+
 \sum_{j=1}^{i-1}(n-\deg P_j)+\epsilon_2\\
 &\qquad\equiv
 \sum_{j=1}^{i-1}(\deg P_j+1)+1\pmod2,
\end{aligned}
\]
这就证明了命题~\ref{prop:8-5-1}(2)。
\end{proof}

于是，命题~\ref{prop:8-5-1} 的 (1) 和 (2) 分别蕴含
\eqref{eq:3-5-17-2} 和 \eqref{eq:3-5-17-1}。因此，我们完成了
$(C(L;\Lambda_{0,\mathrm{nov}}),\mathfrak m)$ 中
$\widehat d\circ\widehat d=0$ 的证明。
